\subsection{Relation transport with the original see-saw scalars}
\label{sec:theta-relation-transport}

The remaining comparison of the actual theta chain with the normalized
Poincare operators can be separated into a family realization and an
exact relation-transfer argument. In this section the family
realization hypotheses are retained. In particular we do not infer a
common real test function for a family from the separate existence of
test functions for its members.

Use the common family of the source \(\mathrm{ValLastX}\):
the length is \(M=2n+2\), the index \(t\) ranges over the original row
tuples, and the intermediate column and row operators are
\(S_{t,j},K_{t,j},\operatorname{io}_{t,j}\), with
\(\operatorname{ioT}_{t,M-2}\) the last row operator.
Let \(\operatorname{per}_{t,j}\) denote the intermediate Bessel
period, with its actual character and translation sums, and let
\(\operatorname{tail}_t\) be the terminal bilinear period.
We use all the original scalars, not just their absolute values:
\[
 A_{t,j}=cV_t(j)cK_t(j)cK_t(j+1),\qquad
 B_{t,j}=cI_t(j)cI_t(j+1),
\]
\[
 C_t=\left(\prod_{i=0}^{n}A_{t,2i}\right)cP_t,\qquad
 D_t=\prod_{j=0}^{M-1}cI_t(j).
 \tag{RT1}
\]
Thus \(C_t=cTop_t(n+1)\), \(D_t=cBot_t\).

\begin{proposition}[Exact downward transport of finite relations]
\label{prop:theta-downward-relations}
Assume the original \(\mathrm{LawsX}\) on this common family,
including its local unfolding laws and the original real
intertwiner spanning assertions. Use the ordinary analytic
two-step unfolding and its density assertion on the prescribed
automorphic spaces. Then every finite relation
\[
 \sum_i a_i C_{t_i}\,
   \operatorname{tail}_{t_i}(\operatorname{up}_{t_i}(n+1)\psi_i,F_i)=0
 \tag{RT2}
\]
implies, for every common pair of real tests \(y,v\),
\[
 \sum_i a_i D_{t_i}\,
 \operatorname{per}_{t_i,0}
   \bigl(\psi_i,
   \operatorname{dn}_{t_i}(M-1)\operatorname{ioT}_{t_i,M-2}F_i;
     y,v\bigr)=0 .
 \tag{RT3}
\]
Here \(\psi_i\) belongs to the original ramification subspace,
\(\operatorname{up}(j+1)=K_{2j+1}K_{2j}S_{2j}\operatorname{up}(j)\)
and \(\operatorname{dn}\) is the ordered composition of all
indicated lower row operators. No scalar has been cancelled on the
assumption that it is a unit.
\end{proposition}

\begin{proof}
We first verify the common finite test coefficients in the source's
double-lift identity. At an even intermediate step \(j\) and for one
of the real bi-Weil test functions \(C\) of \(\mathrm{ValX}\) or
\(\mathrm{ValLastX}\), use the usual real commuting Weil actions
and the finite-degree Fock description of the prescribed central
eigenvectors. The real reduction used here is the source's
\(\mathrm{doubleLift\_reduction\_to\_fockVectors\_uniform}\):
it writes the double lift as a finite sum of consecutive single
lifts with pairs of real polynomial vectors \((P_\nu,P'_\nu)\).
Its uniformity across \(t\) is essential. In applying that
reduction, \(\mathrm{RealEq}\) identifies the two real actions of
every member with the same reference real actions, and
\(\mathrm{RdX}\) supplies their common central character. The
decomposition is of the fixed real bi-test, not of a member's
individual finite-adelic theta kernel. Thus the finite family of
pairs can be used in every member. If the relevant index sets
are empty, all double-lift sums are zero. In this paragraph the
real reduction theorem itself is retained as an analytic input.

For each \(P_\nu\), \(\mathrm{KerX}\) writes the sum of the
single lifts at level \(j\) as
\[
 cK_t(j)\sum_{z}a_\nu(z)
       \operatorname{rd}_{t,j+1}(K_{t,j}\psi,z)
\]
in each of the intermediate finite class indices.
The coefficients \(a_\nu(z)\) have finite support and are
independent of \(t,\psi\) and of the output class index.
Applying the second single-lift law with \(P'_\nu\) to each term
gives the corresponding coefficient family \(b_{\nu,z}(y')\).
Finite-sum linearity of the integral lifts is valid on the
continuous automorphic functions of \(\mathrm{RdX}\).
Therefore the exact common double-lift coefficient is
\[
 a(y')=\sum_{\nu,z}a_\nu(z)b_{\nu,z}(y').
 \tag{RT4}
\]
The resulting two-lift sum has coefficient
\(cK_t(j)cK_t(j+1)\,a(y')\), with the final column
\(K_{t,j+1}K_{t,j}\psi\). This proves the finite-sum part of
\(\mathrm{TwoX}\) with a coefficient family independent of \(t\).
The original ramification conditions are preserved by the two
\(K\)'s, so both applications of \(\mathrm{KerX}\) are legal.

For an ordinary two-step descent, suppose a finite weighted
relation of the upper periods is zero for every real test.
Test the lower periods first on
\(I_j(P_1,I_{j+1}(P_2,v''))\), where the \(I\)'s are the common
real intertwiners. The two lift identities in \(\mathrm{Tup.Is}\)
identify the result, multiplied by \(B_{t,j}\), with the two
consecutive analytic lifts. The finite translation sums are
linear. The actual double-step unfolding expresses the
relevant integral by the iterated integrals of
\(\mathrm{ValX}\). We use its density and absolute-integrability
hypotheses, in particular the local compact-coset support
\(\mathrm{LocFin}\) and rapid decay on the original cuspidal
row spaces; no interchange of merely conditionally convergent
integrals is made.

At each pair of real translates in this unfolding,
\(\mathrm{ValX}\) replaces that iterated integral by \(cV_t(j)\)
times the period with the real bi-test. Apply (RT4) with
\(S_{t,j}\psi\) in place of \(\psi\). The last real translate
of the row vector is a common real linear change \(v''\mapsto
v_*\), by the right-translation identity for the real model.
The resulting finite weighted sum is
\[
 \sum_{y'}a(y')\sum_i a_i A_{t_i,j}\,
   \operatorname{per}_{t_i,j+2}
   (K_{t_i,j+1}K_{t_i,j}S_{t_i,j}\psi_i,G_i;y',v_*).
 \tag{RT5}
\]
Every inner sum is zero by the assumed relation, since the
tests and coefficients in (RT5) are common to the family.
The original unfolding density theorem then gives zero for
the lower relation on the full real spaces. This is
\(\mathrm{StepX}\) with lower coefficient \(a_iB_{t_i,j}\).

At the last pair of steps the same calculation uses
\(\mathrm{ValLastX}\) and the terminal integral in place of the
iterated integrals. Here \(\mathrm{PerX}\) identifies each
terminal integral on a common test \((y',v'')\) with
\(cP_t b(y',v'')\operatorname{tail}_t\), with \(b\) independent
of the row tuple. Thus the lower sum on the intertwiner tests
is
\[
 \sum_{y'}a(y')b(y',v'')\sum_i
   a_i A_{t_i,M-2}cP_{t_i}\,
     \operatorname{tail}_{t_i}
     (K_{t_i,M-1}K_{t_i,M-2}S_{t_i,M-2}\psi_i,F_i).
 \tag{RT6}
\]
Again every inner sum is zero. Linearity is justified on the
original compact terminal quotient; spanning of the common
real intertwiners gives the full lower relation.
This is \(\mathrm{LastX}\).

Finally in (RT2) use \(\mathrm{LastX}\) at \(j=M-2\),
taking coefficient \(a_i\prod_{k<n}A_{t_i,2k}\) and column
\(\operatorname{up}_{t_i}(n)\psi_i\).
The relation at level \(2n\) has coefficient
\(a_i B_{t_i,2n}\prod_{k<n}A_{t_i,2k}\) and the row
\(\operatorname{io}_{t_i,2n}\operatorname{ioT}_{t_i,2n}F_i\).
Descend by \(\mathrm{StepX}\) successively at
\(2n-2,2n-4,\ldots,0\).
At each stage replace one final factor \(A_{t,j}\) by
\(B_{t,j}\) and prepend the two indicated lower row maps.
This is a finite induction, without division by any of the
\(A_{t,j}\).
At level zero the remaining coefficient is precisely
\(a_iD_{t_i}\). All the upward columns have stayed in the
prescribed ramification subspaces by \(\mathrm{stabS}\) and
\(\mathrm{stabK}\). This proves (RT3), also when \(n=0\),
where only the last step is used.
\end{proof}

\paragraph{Recovering the bottom relations and the arithmetic row
coefficients.}
Keep the actual family links \(\mathrm{LinkZero}\),
\(\mathrm{LinkTail}\), \(\mathrm{LinkDown}\), \(\mathrm{LinkOnto}\),
\(\mathrm{LinkComm}\), \(\mathrm{RatioQ}\), and
\(\mathrm{ServedQ}\). Let \(S\) be the similarity on the word's
column space supplied by \(\mathrm{LinkComm}\), with
\(\|S\phi\|=r\|\phi\|\), \(r>0\), and let \(s\ne0\) be its
common tail scalar. Put
\[
 U=S^{-1},\qquad c_t=\frac{C_t s}{D_t}.
 \tag{RT7}
\]
This uses \(D_t\ne0\) from the lift laws and \(C_t\ne0\)
from \(\mathrm{RatioQ}\). In particular \(U\) is a similarity of
ratio \(r^{-1}\). To prove the bottom two-term relation of
\(\mathrm{BotDat.Rel}\), lift its test row \(T_0\) at the base tuple
by \(\mathrm{LinkOnto}\) and its row \(F\) at tuple \(t\) by
\(\mathrm{LinkDown}\). Write the corresponding upper rows as
\(G_0,G_t\).
Use columns \(\operatorname{col}_0(t_0,U\phi_0)\) and
\(\operatorname{col}_0(t,U\phi)\). Define \(p(\phi_0)\) to be
the terminal period of the first column on \(G_0\), multiplied
by \(C_{t_0}/D_{t_0}\), just as in \(\mathrm{RelOfLinks}\).
The tail and commutation laws give the second scalar as
\(c_t\operatorname{perCol}_{u,t}(\phi,F)\).
If
\[
 \alpha\,p(\phi_0)+
 \alpha'\,c_t\operatorname{perCol}_{u,t}(\phi,F)=0 ,
 \tag{RT8}
\]
apply Proposition~\ref{prop:theta-downward-relations} with the
two coefficients \(\alpha/D_{t_0}\), \(\alpha'/D_t\).
At level zero the \(D\)'s cancel \emph{exactly}, and
\(\mathrm{LinkZero}\) identifies the two periods with the
original bottom periods on \(T_0\) and the actual row.
This proves the bottom relation with no unit assumption on
the individual \(cI,cK,cV,cP\).

The arithmetic property of the ratios in (RT7) is especially
simple. If \(w\) is a served prime of the coefficient field \(K\)
and \(t=\mu(w)\) its original served tuple, \(\mathrm{RatioQ}\)
states
\[
 \left(c_t/c_{t_0}\right)^b=q\in\mathbb Q^\times,\qquad b>0,
 \tag{RT9}
\]
with every prime supporting \(q\) among the rational primes
of the auxiliary set of \(t\).
The common \(s\) cancels in this ratio.
By \(\mathrm{ServedQ}\) none of those rational primes is below
\(w\). If \(c_t/c_{t_0}\) is in the image of \(K\) under the
chosen embedding, its preimage \(k\) satisfies \(k^b=q\).
Thus \(b\,v_w(k)=v_w(q)=0\), proving the exact valuation-zero
claim of \(\mathrm{UnitC}\). No bound on the extension generated
by a complex root in (RT9) is used.

\begin{lemma}[Anchoring and descent without adjoining row factors]
\label{lem:theta-anchor-arithmetic-descent}
Suppose the actual bottom relation has been transported to the
Poincare pairings by \(\mathrm{BotRel}\), with its column
isometry \(R\). Suppose a nonzero Poincare witness exists at
the prescribed common real vector. Then the normalized identity
of \(\mathrm{MIstar2}\) gives the same base pairing on the original
rational row spaces, with row coefficients in the original field
\(K\) satisfying the served-prime valuation-zero property.
The nonzero analytic scalar of each word must be retained.
\end{lemma}

\begin{proof}
Use the actual nonzero pairing \(b_0\) at the witness row
\(T_0\) and column. The similarity \(U\) and isometry \(R\)
are invertible in the finite-dimensional column space, so take
its preimage \(\phi_0\) under \(RU\).
For the function \(p\) in the bottom relation, \(p(\phi_0)\ne0\):
otherwise the relation with coefficients \(1,0\) would force
\(b_0=0\) by \(\mathrm{BotRel}\).
Apply the relation to \((T_0,\phi_0)\) and the actual row
\((t,\phi,F)\) using the coefficients
\(c_t\operatorname{perCol}_{u,t}(\phi,F)\) and \(-p(\phi_0)\).
It gives
\[
 c_t\operatorname{perCol}_{u,t}(\phi,F)b_0
       =p(\phi_0)b(t,\phi,F).
 \tag{RT10}
\]
The proportionality is therefore nonzero, not just a conclusion
about kernels. When \(RU\) is replaced by the norm-preserving
map \(rRU\), the removed common scalar is moved into \(b_0\).
The original coefficient conjugation transports this map to
the column-coordinate space. Conjugation does not change its
norm, by the original norm dictionary for those coordinates.
After applying the original row shift (an isometry on the
restricted row space), and the common nonzero Poincare--Bessel
pairing conversion, (RT10) has the form
\[
 z_{u,t}\,\sigma(A_{u,t}(\phi,F))
 =a_u\langle E V\Theta_t(F),\,P_u U_u\phi\rangle,\qquad
 z_{u,t}=\frac{c_t}{c_{t_0}},\quad z_{u,t_0}=1,\quad a_u\ne0 .
 \tag{RT11}
\]
Here \(\sigma:K\hookrightarrow\mathbb C\) is the chosen
embedding, \(A_{u,t}\) is the original algebraic period,
and the coordinates, possible conjugation of the \(c_t\), and
row and column maps are the source's
\(\mathrm{MIid}\) ones. The displayed \(z_{u,t}\) is read after
that same conjugation. We retain the actual row realization
and norm dictionaries of this identity; (RT10) alone does not
construct them.

For clarity, the remaining arithmetic descent in (RT11) can
be done without adjoining any analytic row ratio. Let
\(P_q(t)\) be the prescribed rational row quotient and
\(i_C\) the base row coordinates. At \(t_0\),
\(\Theta_{t_0}(F)=i_C(P_q(t_0)F)\).
Since \(z_{u,t_0}=1\), the base period vanishes on
\(\ker P_q(t_0)\). It therefore defines the original
rational base pairing \(B_u(\phi,y)\) on the quotient:
choose a rational right inverse of the surjection
\(P_q(t_0)\), or equivalently factor through its kernel.
No integral right inverse is asserted or needed.
The original row-coordinate identity \(\mathrm{key\_row}\)
identifies \(\Theta_t(F)\) with \(i_C(M_tP_q(t)F)\),
where \(M_t\) is the original rational row matrix.
Consequently (RT11) gives the equality
\[
 \sigma\bigl(B_u(\phi,M_tP_q(t)F)\bigr)
     =z_{u,t}\sigma\bigl(A_{u,t}(\phi,F)\bigr).
 \tag{RT12}
\]
If \(z_{u,t}\in\sigma(K)\), use its unique preimage
\(\zeta_{u,t}\in K^\times\). Otherwise (RT12) forces the
entire pairing \(A_{u,t}\) to vanish: a single nonzero
rational entry would express \(z_{u,t}\) as a ratio of two
elements of \(\sigma(K)\).
In the second case the left side also vanishes on every
rational input. Set \(\zeta_{u,t}=1\) in that case.
Thus for all the original rational inputs,
\[
 B_u(\phi,M_tP_q(t)F)=\zeta_{u,t}A_{u,t}(\phi,F),\qquad
 \zeta_{u,t_0}=1 .
 \tag{RT13}
\]
At a served \(w\), (RT9) proves \(v_w(\zeta_{u,\mu(w)})=0\)
in the first case; in the second it is immediate. If the
source uses conjugate row factors, apply the transported
valuation condition to those factors, as specified in
\(\mathrm{MIUnit2}\). Thus the correction preserves every
required local divisibility length in the actual row tuple.
We have made no assertion about its valuations at other primes.

For the norm transfer at the base tuple the original
normalization is essential. If \(n_P>0\) is its prescribed
normalizer, then the base coordinate frame satisfies
\[
 \|y\|_{X_E}=n_P^{-1/2}\|i_Cy\|,\qquad
 \|Ee\|=n_P^{-1/2}\|e\|\quad(e\in EEl).
 \tag{RT14}
\]
The row shift \(V\) preserves \(EEl\) and its norm.
Since \(i_C\) identifies the base coordinate space with
\(EEl\), \(W=i_C^{-1}Vi_C\) preserves the frame norm in
(RT14). The column map \(U_u\) likewise preserves the
original \(X_c(u)\) norm. The base case of (RT11),
extended by complex bilinearity (the row anti-linearity of
\(E\) is paired with the conjugate-linear argument of the
Hermitian pairing), is therefore
\[
 \sigma B_u(x,y)
       =a_u\langle E i_C W y,\,P_u U_u x\rangle .
 \tag{RT15}
\]
This is the source's norm factorization \(\mathrm{L1e}\),
including the nonzero scalar \(a_u\).
It gives the arithmetic cut-pairing identity using (RT13)
on the original rational quotient, while keeping the
restricted Hilbert spaces and frame norms unchanged.
In particular we do not replace \(a_u\) by one, or infer
a uniform absolute-height bound for all the \(\zeta_{u,t}\).
\end{proof}

The preceding arguments supply the finite-test, relation-induction,
anchoring and arithmetic descent parts of the comparison.
The outstanding analytic realization is not implicit in them:
one must still construct \(\mathrm{LawsX}\) and
\(\mathrm{FamAllQ}\) from the actual finite-prime chain,
including its real bi-test decomposition, local support,
both unfolding laws, row lift and onto/commutation assertions,
and the exact ratios \(\mathrm{RatioQ}\) on the served tuple
sets. The original nonzero Poincare family and the restricted
norm and level dictionaries are also indispensable.
No complete \(\mathrm{MIstar2}\) or W06 conclusion is asserted
without these suppliers.
